% Propietario conservado: /Users/ruben/Documents/New project/output/REVISION_ARGUMENTAL_APERTURAS_CIERRES_20260915/VI/delivery/MANUSCRIPT_VI_ES_EN_COMPLETE/source_es/sections/07_contraste_metrologico.tex
% Artículo VI. Propietario: indice_articulos. Adición autónoma de exposición.
% Procedencia: RESULTADO_RECUPERADO; controles documentales nuevos separados.
% REV2/fuente/modulos/15_comparacion_cuantitativa_integral.md (completo).
% REV2/fuente/modulos/16_precision_robustez_y_contraste.md (completo).
% REV2/fuente/modulos/19_realizacion_exhaustiva_estado_por_estado.md (completo).
% Integral: integracion_83/masas_p0/corpus/source/masas/
% 19_comparacion_cuantitativa_posterior.tex, parte anterior a \endinput.
% Integral: integracion_83/masas_p0/tex/apendice_inventario_471_384_sucesor.tex.
% Depende de la construcción de atlas, firma, torres y realización material
% incorporada por el editor en las secciones precedentes. No la sustituye.

\section{Realization of observables and metrological comparison}
\label{vi:sec:vi-contraste-metrologico}

The value of a mass acquires its meaning within the structure that
produces it and the representation in which it is measured. The
APP--TRIT--TPK construction first determines the admissible trajectories, their oriented
records and the signature readers; the towers and the material realization
act on those objects. Experimental comparison is placed at the end of
this composition. Its function consists in comparing observables of the same
type, preserving the identity of the state and quantifying agreement in a
common chart. This position allows the mathematical structure,
the evaluated coordinate and its physical interpretation to be studied together.

This section brings together three levels that perform different functions:
the internal domain of routes and their quotients; the classes and evaluations
individualized in the mass owners; and the external inventory of
observables used as a reference. The order of presentation preserves the
constructive direction of the article. External quantities appear as
terms of comparison, while internal values retain the signatures,
operations and conditions that fix their evaluation.

\subsection{Structural domain, material classes and external inventory}
\label{vi:subsec:vi-dominios-contraste}

Let \(\mathcal C\) be the set of APP cells and \(\mathcal R\) its domain
of oriented routes. The identification of a route contains the cell
coordinates and the two initial orientations:
\[
 \gamma=(r,c,h_0,v_0),\qquad
 r,c\in\{1,\ldots,9\},\quad h_0,v_0\in\{-1,+1\}.
\]
Therefore, \(\#\mathcal C=81\) and \(\#\mathcal R=81\cdot4=324\).
The \(56\) trajectory families and the \(13\) multisections belong
to quotients and structural groupings of the same domain, with the
relations and readers established above. These four cardinalities
describe distinct internal objects. The assignment of a physical
representation can preserve an entire fiber; a species label does not
in itself impose a pointwise choice within it.

The public table of \(23\) classes organizes the realizations by sector.
Its function is to relate internal residence to the representation,
charge, spin, color, flavor and type of observable. A class can
comprise several species, and a catalog identity can bring together
different charge states. The correspondence is expressed through typed
maps and preserves these multiplicities.

The frozen metrological inventory contains \(471\) mass records and
\(384\) width records. The corpus documents \(1170\) identities in \(438\)
groups, of which \(403\) contain at least one of those observables.
The sum
\[
 471+384=855
\]
counts external summary-table observables. The atlas \(81/56/13/324\),
the \(23\) classes and the \(855\) records thus belong to three levels
of description. Their association preserves the type of each set, instead
of turning their cardinalities into a single number of particles or predictions.

\begin{center}
\begin{tabular}{p{0.49\linewidth}rrr}
\toprule
Catalog family & Groups & Masses & Widths\\
\midrule
General & 46 & 53 & 9\\
Quark sector & 6 & 8 & 1\\
Mesons & 201 & 243 & 225\\
Baryons & 149 & 166 & 149\\
External gravitational record & 1 & 1 & 0\\
\midrule
Total with observable & 403 & 471 & 384\\
\bottomrule
\end{tabular}
\end{center}

The gravitational row preserves the identity appearing in the external source;
its presence in the inventory does not define an elementary particle of HMT.
The remaining \(35\) groups retain their catalog identity, although
the snapshot does not assign them a summary-table mass or width. Documentary
integrity encompasses both facts: the quantities recorded and the identities
for which that snapshot does not publish one of those quantities.

\subsection{The observable before the numerical coordinate}
\label{vi:subsec:vi-tipos-observable}

A material output is specified through the state, the representation,
the operator and the reading procedure. For a fiber \(F_X\), let
\(\widehat M_X\) be the mass operator already constructed and let \(P_{X,a}\)
be the projector corresponding to the state or preparation channel \(a\).
For \(\Psi_X\) normalized in its image, the spectral distribution is
\[
 \mu_{X,a}(B)
 =\langle\Psi_X,
   P_{X,a}E_{\widehat M_X}(B)P_{X,a}\Psi_X\rangle,
\]
where \(E_{\widehat M_X}\) is the spectral measure of the operator. This
expression preserves the prepared state and the multiplicity of the fiber.
A single scalar coordinate for every state prepared in that subspace
is obtained under the condition
\[
 \operatorname{Ran}(P_{X,a})\subset\operatorname{Dom}(\widehat M_X),
 \qquad \widehat M_XP_{X,a}=m_XP_{X,a}.
\]
Indeed, each vector of that subspace is then an eigenvector of the
mass operator for the same value \(m_X\), and its spectral measure is
concentrated at that value. If the compression equality
\(P_{X,a}\widehat M_XP_{X,a}=m_XP_{X,a}\) is used, it is additionally required that
\(P_{X,a}\) reduce \(\widehat M_X\); in finite dimension, this is equivalent
to their commutation. Without that condition, compression may fix only
the mean value. When the spectral measure retains several values, the
output is presented as a distribution or multiplet with its weights.

For a fixed density \(\rho\) in a finite fiber, let \(m_j\) be
the distinct eigenvalues and \(E_j\) their spectral projectors.
The weights \(p_j=\operatorname{tr}(\rho E_j)\geq0\), with
\(\sum_jp_j=1\), determine
\[
 \overline m=\sum_jp_jm_j,\qquad
 \operatorname{Var}_{\mu}(m)=\sum_jp_j(m_j-\overline m)^2.
\]
\paragraph{Scalar concentration criterion for the fixed state.}
The variance vanishes if and only if all values with positive weight
coincide. Indeed, each term is nonnegative; the sum is zero
precisely when \(m_j=\overline m\) for every \(j\) with
\(p_j>0\). This determines when a mass distribution admits
a scalar publication without loss of the dispersion of the prepared state.
This criterion refers to the fixed density and differs from the operator
condition on the entire subspace. The electronic case uses its
operator on the central fiber, as previously constructed;
the rank one of an arbitrary projector is not sufficient to concentrate the spectral
measure of another operator.

Mass and width use coordinated readers of different types. For a
resonant state, a pole convention can be expressed in energy
coordinates as
\[
 z_X=M_Xc_{\rm int}^{2}-\frac{i}{2}\Gamma_X.
\]
The real part is the mass-energy coordinate; \(\Gamma_X\) has
the dimension of energy and measures the width. Once the temporal chart is fixed,
\(\tau_X=\hbar_{\rm int}/\Gamma_X\). The relation assumes the pole and
survival regime for which it was constructed; a rate, a width and
a lifetime retain their units and their conversion rules. This
distinction explains why the inventory preserves two classes of observable.

The same principle determines comparison by sector. Current quark masses
are expressed with scheme \(\mathscr S\) and scale \(\mu\);
resonances distinguish pole mass, Breit--Wigner parametrization and
on-shell convention. Neutrinos distinguish mass eigenvalues,
squared differences, flavors and mixing parameters. One-sided
bounds retain the documented confidence level and statistical
convention. An upper bound is not converted into a positive central value,
nor a squared difference into an absolute mass.

\subsection{Fixing the construction and correspondence of records}
\label{vi:subsec:vi-congelacion}

The realization of a species \(X\) transmits the necessary information
from its internal antecedents. For comparison purposes, we can gather it
in the record
\[
 \mathfrak D_X=
 \bigl(F_X,\rho_X,P_{X,a},\mathcal L_X,
       \sigma_X,\eta_X,\widehat M_X,\mathcal U_X\bigr),
\]
where \(\rho_X\) identifies the representation, \(\mathcal L_X\) the
genealogical record, \(\sigma_X\in\mathbb Z^5\) the signature,
\(\eta_X\) the tower section and \(\mathcal U_X\) the units
chart. All these data receive their meaning from the
APP--TRIT--TPK operations and the realizations proved in the article. Gathering them
in a tuple is a documentation operation: their constructions and
proofs remain in the sections that produce them.

The external reference is represented separately by
\[
 o=(\mathrm{id},\mathrm{especie},\mathrm{observable},
      v,\mathrm{unidad},\mathscr S,\mu,
      \mathrm{incertidumbre},\text{limit},\mathrm{confianza},
      \mathrm{fuente},\text{edition}).
\]
The association \((\mathfrak D_X,o)\) must be established through
physical identity and the type of observable. In an extensive table, this
operation is performed through identity and provenance keys; the proximity
of two figures does not serve as a key. The record retains the
particle or antiparticle label, the group that possesses the observable and
the identity of the measured property.

The independence of the selection procedure is examined by hiding or
permuting the metrological column and checking that \(\mathfrak D_X\)
remains unchanged. That test makes sense for a generator whose
effective composition can be executed; merely preserving an output
file does not in itself constitute such a test. The cryptographic fingerprint
fixes which artifact is compared and makes modifications detectable, but
causal independence belongs to the operation and its dependencies.

Statuses are assigned to the domain and operation that underpin each
statement. A row can simultaneously document an internal operator,
an evaluation given a signature and an external comparison, without identifying
those three levels with each other. Every subsequent realization is linked to the
definition, proposition or table that establishes its scope. In particular, the
catalog's documentary classifications apply to the snapshot specified by
each record and do not modify the mathematical status of a later
construction.

\subsection{Metrological correspondence of the twenty-three classes}
\label{vi:subsec:vi-clases-metrologia}

The following table preserves the \(23\) classes of the source and specifies the
form of reading that each requires upon reaching the metrological interface. Its
function differs from the operator classification of
proposition~\ref{vi:vi:prop:producto-fibrado-especie-ruta} and the documentary
inventory of tables~\ref{vi:tab:vi-estados-clases} and~\ref{vi:tab:vi-clases}:
here the comparable quantity, the realization convention and
the data that must accompany it are distinguished. Flavor and sector names
guide the correspondence; the transmitted mathematical object remains in
the corresponding fibers, representations and readers.

\begingroup
\small
\setlength{\tabcolsep}{3pt}
\begin{longtable}{p{0.20\linewidth}p{0.25\linewidth}p{0.47\linewidth}}
\caption{Metrological reading of the twenty-three classes and type of object transmitted.}
\label{vi:tab:vi-clases-metrologia}\\
\toprule
Class & Representatives & Object documented for the reading\\
\midrule
\endfirsthead
\toprule
Class & Representatives & Object documented for the reading\\
\midrule
\endhead
\bottomrule
\endfoot
Charged lepton 1 & Electron and positron & Central route, conjugation and beta coordinate.\\
Charged lepton 2 & Muon and antimuon & Fiber, readers and evaluation of the declared signature.\\
Charged lepton 3 & Tau and antitau & Fiber, readers and evaluation of the declared signature.\\
Neutrino 1 & Electronic flavor and conjugate & Neutral fiber and mass--flavor realization.\\
Neutrino 2 & Muonic flavor and conjugate & Neutral fiber and mass--flavor realization.\\
Neutrino 3 & Tauonic flavor and conjugate & Neutral fiber and mass--flavor realization.\\
Quark arriba 1 & \(u,\bar u\) & Flavor and color fiber; scheme and scale chart.\\
Quark abajo 1 & \(d,\bar d\) & Flavor and color fiber; scheme and scale chart.\\
Quark arriba 2 & \(c,\bar c\) & Flavor and color fiber; scheme and scale chart.\\
Quark abajo 2 & \(s,\bar s\) & Flavor and color fiber; scheme and scale chart.\\
Quark arriba 3 & \(t,\bar t\) & Flavor and color fiber; scheme and scale chart.\\
Quark abajo 3 & \(b,\bar b\) & Flavor and color fiber; scheme and scale chart.\\
Electromagnetic gauge & Photon & Null chart and helicity representation.\\
Chromatic gauge & Gluons & Null chart and adjoint color representation.\\
Calibre cargado & \(W^+,W^-\) & Route, signature and pole convention of the realization.\\
Calibre neutro & \(Z^0\) & Route, neutral mixing and pole convention.\\
Scalar & Higgs & Scalar representation, signature and pole convention.\\
Mesonic composite & Mesons and conjugates & Constituents, link record and composite reading.\\
Baryonic composite & Baryons and conjugates & Trilinear composition, link and composite reading.\\
Fibonacci topological & \(1,\tau\) & Fusion ring and Frobenius--Perron dimension.\\
Ising topological & \(1,\psi,\sigma\) & Fusion rules and sectoral representation.\\
Cyclic topological & \(\mathbb Z/6\mathbb Z\) & Winding characters and braid representation.\\
Virtual & Secciones intermedias & Finite horizon, prolongation and obstruction.\\
\end{longtable}
\endgroup

In charged leptons, the reading leads to the mass of the central section
or to the evaluation of the selected fiber. In neutrinos, it additionally preserves
the transformation between mass and flavor and the squared differences; therefore
it is not reduced to an isolated scalar. In quarks, every figure is accompanied by
the renormalization scheme and scale. The electromagnetic and
chromatic classes publish, respectively, an experimental bound and the null gauge
section; the \(W\), \(Z\) and scalar classes require the declared pole
or resonance convention. Composites preserve mass, width and constituent
identity. The three topological classes are compared through their
fusion, braiding and gap laws once the dynamical realization is fixed;
their dimensions and characters belong to their own codomains. The virtual
class preserves the pole, the obstruction and the prolongation horizon. This
reading determines which object reaches metrology without confusing the
twenty-three classes with twenty-three universal mass scalars.

\subsection{Metrological snapshot, units and transcription accuracy}
\label{vi:subsec:vi-corte-metrologico}

The snapshot identified by the corpus as Particle Data Group
2026 is used, with a cutoff date of January \(15\), \(2026\) and a documented
consultation on July \(26\), \(2026\). This section reproduces that
documentary state: it does not incorporate a subsequent catalog update.
The references of the evaluation supplement and those of the historical grid
retain their individual provenance; they do not automatically receive the edition
of the general inventory.

Each value is preserved as its original decimal string, together with its
presentation, unit, uncertainty and provenance key. This decision
prevents an intermediate floating-point conversion from changing the last
digits or removing significant zeros from the transcription. The original
string and the typographical representation are two distinct data items: the
former allows the file's arithmetic to be reproduced exactly; the
latter communicates the result with a precision appropriate for reading.
The digits of the digital representation of a reference do not increase its
experimental precision.

Conversion between MeV and GeV uses the exact factor \(10^3\), applied
to the value and its uncertainty. To present a mass, the unit is written
MeV/\(c^2\) or GeV/\(c^2\); to present the rest energy,
MeV or GeV. The relation \(E=mc^2\) is the chart that allows passage from one
reading to the other. A comparison of coordinates requires fixing that chart in
both columns.

The inventory snapshot expressly retains the scheme and scale of its external
records as not supplied. The name of a property
may contain an indication such as ``Breit--Wigner''; it is preserved literally,
but that partial indication does not complete the metadata the comparison
requires. Empty cells, unavailable-data markers and
not-applicable markers remain differentiated. A documentary unknown
is not replaced by zero.

\subsection{Individualized evaluations and reading conditions}
\label{vi:subsec:vi-evaluaciones-documentadas}

Table~\ref{vi:tab:vi-evaluaciones-veinte} brings together \(20\) numerical
presentations: the electronic beta coordinate, two null outputs and
\(17\) evaluations conditional on declared descriptors. Among
the latter, \(11\) correspond to elementary sectors and \(6\) to
composites. This count expresses the contents of that documentary presentation;
the fibers, families and records of the domain retain their own extent.

The table retains the twenty entries. To facilitate reading,
four modalities are identified: \(\mathrm{B}\), electronic beta
coordinate; \(\mathrm{N}\), null chart; \(\mathrm{F}\), evaluation
conditional on a complete signature; \(\mathrm{A}\), conditional angular
evaluation. The subscript \(c\) indicates material composition. The
condition of each entry belongs to its type of result and is also
preserved in the electronic catalog accompanying the text.

\begingroup
\small
\setlength{\tabcolsep}{3pt}
\begin{longtable}{p{0.13\linewidth}p{0.32\linewidth}p{0.38\linewidth}p{0.07\linewidth}}
\caption{Twenty individualized evaluations, with their chart and comparison modality.}
\label{vi:tab:vi-evaluaciones-veinte}\\
\toprule
State & HMT coordinate, MeV/\(c^2\) & Subsequent documentary reference & Type\\
\midrule
\endfirsthead
\toprule
State & HMT coordinate, MeV/\(c^2\) & Subsequent documentary reference & Type\\
\midrule
\endhead
\bottomrule
\endfoot
\(e^\pm\) & \(0.510998950\allowbreak690000809\) & \(0.51099895069(16)\) MeV/\(c^2\) & B\\
\(\mu^\pm\) & \(105.658360294\allowbreak435829303\) & \(105.6583755\pm0.0000023\) MeV/\(c^2\) & F\\
\(\tau^\pm\) & \(1776.860917631\allowbreak432295595\) & \(1776.93\pm0.09\) MeV/\(c^2\) & F\\
\(u,\bar u\) & \(2.165924536\allowbreak173907663\) & \(2.16\pm0.07\) MeV/\(c^2\) & A\\
\(d,\bar d\) & \(4.679550162\allowbreak948874172\) & \(4.70\pm0.07\) MeV/\(c^2\) & A\\
\(c,\bar c\) & \(1270.500838641\allowbreak911135286\) & \(1.2729\pm0.0045\) GeV/\(c^2\) & A\\
\(s,\bar s\) & \(92.940093907\allowbreak845640641\) & \(92.9\pm0.7\) MeV/\(c^2\) & A\\
\(t,\bar t\) & \(172597.479544019\allowbreak136261367\) & \(172.60\pm0.27\) GeV/\(c^2\) & A\\
\(b,\bar b\) & \(4190.119763299\allowbreak790218075\) & \(4.186\pm0.006\) GeV/\(c^2\) & A\\
\(\gamma\) & \(0\) & \(<10^{-18}\) eV/\(c^2\), retained bound & N\\
\(g\) & \(0\) & Null gauge chart; no experimental comparison row & N\\
\(W^\pm\) & \(80378.964909704\allowbreak547024865\) & \(80.3625\pm0.0077\) GeV/\(c^2\) & F\\
\(Z^0\) & \(91187.533444313\allowbreak407844855\) & \(91.1879\pm0.0020\) GeV/\(c^2\) & F\\
\(H\) & \(125292.466042536\allowbreak133000433\) & \(125.13\pm0.11\) GeV/\(c^2\) & A\\
\(\pi^0\) & \(134.976724327\allowbreak872125739\) & \(134.976827767685\) MeV/\(c^2\), supplement & \(\mathrm F_c\)\\
\(\pi^\pm\) & \(139.570485171\allowbreak572191375\) & \(139.570390983681\) MeV/\(c^2\), supplement & \(\mathrm F_c\)\\
\(K^\pm\) & \(493.677085649\allowbreak530612476\) & \(493.676599458041\) MeV/\(c^2\), supplement & \(\mathrm F_c\)\\
\(K^0\) & \(497.611186497\allowbreak750457359\) & \(497.610767531258\) MeV/\(c^2\), supplement & \(\mathrm F_c\)\\
\(p,\bar p\) & \(938.271751546\allowbreak984898299\) & \(938.27208943\) MeV/\(c^2\), supplement & \(\mathrm F_c\)\\
\(n,\bar n\) & \(939.565810472\allowbreak507023313\) & \(939.56542194\) MeV/\(c^2\), supplement & \(\mathrm F_c\)\\
\end{longtable}
\endgroup

In this table, the reference values are presented with the decimal
notation and units declared for each observable. The calculation
record also preserves the decimal string used in the evaluation,
including the difference between that string and its rounded presentation for
the tau, \(Z\) and Higgs. The quark comparison
remains informative until the scheme and scale chart is assembled.
The \(W,Z,H\) sectors also require preservation of the resonance
convention. The six composite references come from the supplement
identified in the evaluation records; their metadata are not completed
by analogy with the elementary references.

The complete signatures documented in modalities \(\mathrm F\) and
\(\mathrm F_c\), in the order
\(\sigma=(n_A,s_C,k_{\Delta_4},b_{120},b_\zeta)\), are:
\[
\begin{array}{c|r@{,\,}r@{,\,}r@{,\,}r@{,\,}r}
\mu&42&-3&8&0&2\\
\tau&64&0&25&-1&6\\
W&94&-1&0&-2&4\\
Z&95&-1&-11&0&-4\\
\pi^0&44&-4&-25&1&-2\\
\pi^\pm&44&1&-2&2&-1\\
K^\pm&54&-1&17&-2&2\\
K^0&54&0&32&3&-2\\
p&59&0&9&1&0\\
n&59&1&-34&0&1
\end{array}
\]
The angular modalities preserve, respectively for
\(u,d,c,s,t,b,H\), the pairs
\[
(n_A,s_C)=(11,7),(17,8),(61,8),(41,-3),(100,-1),(71,-5),(97,9).
\]
The two declared components are not extended with
three implicit zeros. The evaluation equation and the mechanism that
produces a descriptor are different statements; their proofs
are found in their respective definitions and propositions.

\subsection{Relative differences and uncertainty propagation}
\label{vi:subsec:vi-residuales}

Given a common chart and two positive coordinates of the same observable,
the relative difference and the logarithmic difference are
\[
 \delta_X=\frac{m_X^{\rm HMT}-m_X^{\rm ext}}{m_X^{\rm ext}},
 \qquad r_X=\log\frac{m_X^{\rm HMT}}{m_X^{\rm ext}},
 \qquad \delta_X^{\rm ppm}=10^6\delta_X.
\]
A common unit conversion \(m\mapsto am\), with \(a>0\),
leaves both differences invariant: the factor \(a\) cancels in the
quotients. The absolute difference is multiplied by \(a\), as is
its uncertainty. For a zero external value, a one-sided bound or an
interval without a central estimate, the preceding quotient is replaced by
the comparison corresponding to that type of observation.

When a joint uncertainty model is available, let
\(d_X=m_X^{\rm HMT}-m_X^{\rm ext}\). Its variance is
\[
 u^2(d_X)=u^2(m_X^{\rm HMT})+u^2(m_X^{\rm ext})
          -2\operatorname{Cov}(m_X^{\rm HMT},m_X^{\rm ext}).
\]
The formula results from expanding
\(\operatorname{Var}(U-V)\). If \(u(d_X)>0\), the normalized
difference is \(z_X=d_X/u(d_X)\). Its statistical interpretation requires
the corresponding probabilistic model and covariance; an unrecorded field
does not amount to zero variance. Asymmetric uncertainties
are retained as such and are not automatically converted into a symmetric
deviation.

For several species, the joint covariance matrix is retained. With
the discrete signatures fixed, if \(f_X(\vartheta)=\log(m_X/m_{\rm ref})\)
depends on continuous coordinates already constructed and \(J\) is its
Jacobian matrix, first-order propagation gives
\[
 \Sigma_f=J\Sigma_\vartheta J^{\mathsf T}.
\]
This formula preserves the correlations shared by the angular pair,
\(\Delta_4\), the action ratio and the towers. Its application presupposes
the regime in which linearization controls the remaining terms.
An integer change of signature is recorded as a discrete variation of the
construction, not as an infinitesimal Gaussian fluctuation.

The precision of an evaluation therefore has several sources: the
experimental reference; the units, scheme and scale chart; the
distribution of routes; the tower approximation; and the selected
coupling. If \(\Sigma_{ij}\) denotes the covariance between the
contributions \(i,j\), the total covariance of their sum is
\(\sum_{i,j}\Sigma_{ij}\). Keeping the cross blocks avoids
attributing independence to several coincidences that use the same
antecedents.

\subsection{Absolute scale and control of the harmonic approximation}
\label{vi:subsec:vi-precision-torres}

The absolute scale and the relative correction enter at different
places in the construction. Given an internally constructed reference basal
section \(m_\star^{(0)}\), let
\(\chi_\star=\chi_{\rm full}(\sigma_\star,\eta_\star)\).
In the normalization relative to that section, the expression is
\[
 m_X^\beta=m_\star^{(0)}\,
 \frac{\chi_{\rm full}(\sigma_X,\eta_X)}{\chi_\star}
 R_{\rm act}^{\kappa_X}.
\]
The common basal factor remains explicit. If the reference is electronic,
\(m_\star^{(0)}=m_e^{(0)}\) is the dimensional realization of
\(\mathfrak e_0\), with all the factors of its basal equation, and
\(\chi_\star=\chi_{\rm full}(\sigma_e,\eta_e)\). Thus we recover
\(m_e^\beta=m_e^{(0)}R_{\rm act}^{\kappa_e}\). If the common clock is used
to write \(m_\star^{(0)}=m_{\rm clk}b_\star\), the
factor \(b_\star\) preserves that basal normalization; it is not replaced
by \(1\) or incorporated twice into the character. For two states
in the same chart and with the same reference, the quotient is
\[
 \frac{m_Y^\beta}{m_X^\beta}
 =\frac{\chi_{\rm full}(\sigma_Y,\eta_Y)}
        {\chi_{\rm full}(\sigma_X,\eta_X)}
     R_{\rm act}^{\kappa_Y-\kappa_X}.
\]
The cancellation of \(m_\star^{(0)}/\chi_\star\) is exact. It includes
the common basal prefactor and the common factors of the absolute branch of
action, clock and velocity chart. Consequently,
the action decade \(10^{-34}\) retains its role in
the dimensional realization and is not inserted again into the relative
corrector. Differences between routes remain in their signatures, readers
and tower sections.

To control a truncation, it is useful to fix the normal form of the character.
Let
\[
 \zeta_{270}=135\Delta_4^2,\qquad
 N_h=\eta_h+b_{120}\mathbf 1_{\{h=120\}},\qquad
 \mathcal S=\langle90,120\rangle\setminus\{0\}.
\]
Then
\[
\begin{split}
 \log\chi_{\rm full}
 ={}&n_AA_{\rm rad}+\frac{s_C}{6}C^*_{\rm rad}
       +k_{\Delta_4}\Delta_4+b_\zeta\zeta_{270}
       +\sum_{h\in\mathcal S}N_hs_h .
\end{split}
\]
In this convention, the height term \(120\) appears only once.
The term \(b_\zeta\zeta_{270}\) retains its quadratic origin in the
record; \(N_{270}s_{270}\) belongs to the harmonic reading. The
equality of the numerical index does not identify the two objects.

\paragraph{Truncation bound and multiplicative error.}
Suppose that, for \(h>H\), the already determined coefficients of a
route satisfy \(|N_hs_h|\leq Cq^h\), with \(C\geq0\) and
\(0<q<1\). The logarithmic tail \(R_H\) satisfies
\[
 |R_H|\leq\sum_{\substack{h\in\mathcal S\\h>H}}Cq^h
 \leq\sum_{h=H+1}^{\infty}Cq^h
 =\frac{Cq^{H+1}}{1-q}=:\varepsilon_H.
\]
If \(m_H\) is the coordinate obtained by truncating that tower and the other
factors are kept fixed, \(m/m_H=e^{R_H}\). Monotonicity of the
exponential yields
\[
 e^{-\varepsilon_H}\leq\frac{m}{m_H}\leq e^{\varepsilon_H},
 \qquad
 \left|\frac{m-m_H}{m_H}\right|\leq e^{\varepsilon_H}-1.
\]
The bound allows precision to be assigned to the evaluation from the route and its
operators. The values of \(C,q,H\) are justified through that
construction; the experimental residual does not select the length of the
tower or the bound to be proved.

% REMATE_LOCAL_VI_CONTRASTE
\Needspace{29\baselineskip}
\subsection{Complete preservation and reproducibility of the comparison}
\label{vi:subsec:vi-catalogo-reproducible}

The catalog preserves the complete records in tabular and
structured formats. Each row retains its source file, its ordinal,
the original key and the occurrence number of that key. This
identification preserves even two literally identical rows: the
equality of their contents does not eliminate the second occurrence. The change
of format preserves the original strings, empty cells and
statuses, and attaches the fingerprint of the input bytes.

The inventory of \(855\) observables is also linked to the metrological
rows of the integral corpus. The correspondence is verified
by the pair consisting of the observable type and its identifier; the
check compares multiplicities, not just sets of keys. The internal
tables of cells, families, multisections and routes remain separate
from that inventory. The \(23\) classes, the \(20\) numerical
presentations and the historical grid also retain their own files and
statuses.

Transcription checks verify recovery of all
fields, row order, duplicates, decimal strings and
correspondence with the reference tables. Their scope is documentary.
The mathematical proofs of the atlas, the readers and the material realization
are what establish the internal outputs; the metrological comparison
completes that presentation through well-defined observables and charts.
Documentary integrity and the sufficiency of a proof are
checked separately.

Extension of the catalog therefore follows a cumulative rule:
a new evaluation incorporates the state, its record and signature, the tower
section, the reader, the realization, the comparison chart and the
provenance. Previous cases retain their content and their cutoff
date. Thus, growth of the inventory makes a larger part
of the domain visible without retrospectively altering what each proof and
each evaluation had established.
